% Bagian: counterfactuals
% Bab: minimal-change-semantics
% Subbagian: introduction

\documentclass[../../../include/open-logic-section]{subfiles}

\begin{document}

\olfileid[id]{cnt}{min}{int}

\olsection{Pendahuluan}

Stalnaker dan Lewis mengajukan penjelasan tentang kondisional kontrafaktual
seperti ``seandainya korek api itu digesek, korek itu akan menyala.''
Penjelasan mereka merupakan usulan tentang cara memahami secara tepat syarat
kebenaran kalimat semacam itu. Gagasan di balik kedua usulan tersebut adalah
sebagai berikut: untuk menilai apakah suatu kondisional kontrafaktual benar,
kita harus mempertimbangkan dunia-dunia mungkin yang perbedaannya dari keadaan
dunia aktual seminimal mungkin sehingga antesedennya menjadi benar. Jika
konsekuennya benar di dunia-dunia mungkin ini, kontrafaktual tersebut benar.
Sebagai contoh, andaikan saya memegang sebatang korek api dan sekotak korek api.
Di dunia aktual, saya hanya memandang keduanya dan memikirkan apa yang akan
terjadi seandainya saya menggesek korek api itu. Perubahan minimal dari dunia
aktual tempat saya menggesek korek api ialah dunia tempat saya memutuskan untuk
bertindak dan menggeseknya. Perubahan itu minimal dalam arti bahwa tidak ada
hal lain yang berubah: saya tidak sekaligus melompat ke udara, menggesek korek
api tidak sekaligus membakar rambut saya, jari-jari saya tidak tiba-tiba
kehilangan seluruh kekuatannya, saya tidak pada saat yang sama diguyur air
dalam penyergapan dengan SuperSoaker, dan sebagainya. Dalam kemungkinan
alternatif itu, korek api tersebut menyala. Jadi, benar bahwa seandainya saya
menggesek korek api itu, korek tersebut akan menyala.

Penjelasan intuitif ini dapat dipadukan dengan semantik formal bagi logika
kontrafaktual. Lewis memperkenalkan simbol ``$\boxright$'' untuk kondisional
kontrafaktual, sedangkan Stalnaker memakai simbol~``$>$''. Kita akan
memakai~$\cif$ dan menambahkannya sebagai konektif biner ke dalam logika
proposisi. Jadi, selain !!{formula}s berbentuk $!A \lif !B$, kita juga
mempunyai !!{formula}s berbentuk~$!A \cif !B$. Seperti semantik relasional
bagi logika modal, semantik formal ini didasarkan pada model-model tempat
!!{formula}s dinilai pada dunia-dunia, dan syarat keterpenuhan yang
mendefinisikan $\mSat{M}{!A \cif !B}[w]$ diberikan dalam hubungannya dengan
$\mSat{M}{!A}[w']$ dan $\mSat{M}{!B}[w']$ untuk suatu dunia (lain)~$w'$.
Dunia $w'$ yang mana? Secara intuitif, dunia yang paling dekat dengan~$w$ di
antara dunia-dunia tempat~$\mSat{M}{!A}[w']$. Karena itu, suatu relasi
``kedekatan'' juga harus disertakan dalam model.

Lewis memperkenalkan cara yang bermanfaat untuk merepresentasikan situasi
kontrafaktual secara grafis. Setiap dunia mungkin berada di pusat suatu
himpunan bola bersarang yang memuat dunia-dunia lain---bola-bola ini kita
gambar sebagai lingkaran sepusat. Dunia-dunia yang berada di antara dua bola
semuanya sama dekat dengan dunia di pusat; dunia-dunia yang termuat
dalam bola yang lebih dalam lebih dekat, sedangkan dunia-dunia dalam bola yang
mengelilinginya lebih jauh.
\begin{center}
\begin{tikzpicture}[scale=.7]\small
  \spheresystem{5}
  \draw (0,0) node {$w$};
  \propositionintersect{0}{4}{40}{3.5}
  {\tikzset{proposition/.append={smooth,tension=1.4}}
  \proposition[shift={(1.6,-1)}]{90}{1}{30}{4}}
  \path (1.5,3) node[below] {$\formula{B}$};
  \path (3,0) node {$\formula{A}$};
\end{tikzpicture}
\end{center}
Dunia-$!A$ yang paling dekat adalah dunia-dunia $w'$ tempat~$!A$ dipenuhi
yang terletak dalam bola terkecil di sekitar dunia pusat~$w$ (daerah abu-abu).
Secara intuitif, $!A \cif !B$ dipenuhi pada~$w$ jika $!B$ benar pada semua
dunia-$!A$ yang paling dekat.

\end{document}
